\begin{table}[htbp]\centering\small
\caption{Lo proyectado para 2027 en abril y en septiembre de 2026}\label{cua:precrit}
\begin{tabular}{lrrr}
\toprule
Supuesto para 2027 & Pre-Criterios (6 abr.) & CGPE (8 sep.) & Diferencia \\
\midrule
PIB nominal 2027 (mmp) & 39.732,4 & 39.419,4 & $-$313,0 \\
Deflactor del PIB (\%) & 4,0 & 4,0 & 0,0 \\
Inflación INPC dic./dic. (\%) & 3,0 & 3,0 & 0,0 \\
Tipo de cambio promedio & 18,5 & 17,9 & $-$0,6 \\
Tipo de cambio fin de periodo & 18,6 & 18,0 & $-$0,6 \\
Cetes 28 días, promedio (\%) & 5,9 & 6,1 & 0,2 \\
Mezcla mexicana (dólares por barril) & 54,7 & 61,8 & 7,1 \\
Plataforma de producción (mbd) & 1.806,1 & 1.800,1 & $-$6,0 \\
Plataforma de exportación (mbd) & 427,6 & 426,6 & $-$1,0 \\
Crecimiento real (rango, \%) & [1,9\,--\,2,9] & [1,5\,--\,2,5] & $-$0,4 en los dos extremos \\
\bottomrule
\end{tabular}
\\[2pt]\begin{minipage}{0.92\textwidth}\scriptsize Fuentes: Pre-Criterios 2027, Anexo I (p. 34), publicados el 6 de abril de 2026; y CGPE 2027, edición SHCP, Anexo II.5 (p. 66). \textbf{Los dos son estimaciones del mismo Ejecutivo sobre el mismo ejercicio, separadas por cinco meses.} Ninguna de estas cifras admite variación real: son precios, tasas, volúmenes y un nivel nominal de PIB.\end{minipage}
\end{table}
